\documentclass[10pt]{article}
\usepackage[a4paper,margin=0.55in]{geometry}
\usepackage{amsmath,amssymb,mathtools}
\usepackage[dvipsnames]{xcolor}
\usepackage{booktabs}
\usepackage[colorlinks=true,linkcolor=blue]{hyperref}
\numberwithin{equation}{section}
\setlength{\emergencystretch}{3em}\sloppy
\allowdisplaybreaks
\newcommand{\MSb}{\overline{MS}}
\newcommand{\LO}{{\rm LO}}
\newcommand{\new}[1]{{\color{blue}#1}}
\newcommand{\old}[1]{{\color{red}#1}}
\title{Group II Row IV, step by step:\\ where $\log\frac{\vee}{k^+}H_4$ comes from, and what is wrong with the $d^3N^y$ term of Eq.~(5.12)}
\date{}
\begin{document}
\maketitle

\section{The short answer}
\begin{itemize}
\item \textbf{Your advisor is right.} The $d^3N^y_{\xi\rm LO}$ term of Eq.~(5.12), \emph{as written}, has an intrinsic problem:
\begin{itemize}
\item the log stands outside the integral, and then the term is identically zero for $\mathbf k\neq0$;
\item it is not manifestly symmetric under amplitude $\leftrightarrow$ conjugate ($y\leftrightarrow y'$, $x\leftrightarrow x'$);
\item ``LO at $y'=y$'' is not a cross section.
\end{itemize}
\item \textbf{You are also right.} BFKL$\otimes$LO needs $+\log\frac{\vee}{k^+}H_4$, and Row IV contains it. It sits inside the DGLAP integral, at its lower end $\xi\to\frac{k^+}{\vee}$ (Sec.~\ref{sec:small}).
\item \textbf{How the two fit together.} When $+\log\frac{\vee}{k^+}H_4$ is taken out of Row IV into the BFKL list, the same amount must be subtracted from the DGLAP integral. That subtraction is what the $d^3N^y$ term was meant to be. Written correctly it is
\begin{equation}
\boxed{\;-2N_c\frac{1-\xi}{\xi}\frac{d^3N^y_{\xi\rm LO}}{d^3k}\Big[\log\frac{X^2\mu^2}{4}\Big]\;\longrightarrow\;-2N_c\frac{1-\xi}{\xi}\Big(-\frac{4\pi\xi^2}{\mathbf k^2}\Big)\frac{d^3N^y_{\rm LO}}{d^3k}\;}\label{eq:repl}
\end{equation}
with $\frac{d^3N^y_{\rm LO}}{d^3k}$ symmetric (both gluon points at the same $z$, no phase), Eq.~\eqref{eq:LOy}. After the $\xi$ integral it is $-\big[\log\frac{\vee}{k^+}-1+\frac{k^+}{\vee}\big]H_4$.
\item \textbf{The advisor's ``it should match or cancel with $H_4$''} is exactly true:
\begin{itemize}
\item the subtraction is proportional to $H_4$;
\item together with the $-(1-\frac{k^+}{\vee})H_4$ line it gives $-\log\frac{\vee}{k^+}H_4$;
\item together with the BFKL piece $+\log\frac{\vee}{k^+}H_4$ it gives zero, so the full Row IV has no explicit $H_4$ (Sec.~\ref{sec:book}).
\end{itemize}
\item \textbf{The running coupling} is $\LO|_{\alpha_s\to\alpha_s(X^2)}$, $X=y'-y$, exactly your advisor's $\alpha_s(X^2)=\alpha_s+\frac{\alpha_s^2\beta_0}{4\pi}\ln(X^2\hat\mu^2)$ (Step 4).
\end{itemize}
Notation: $X\equiv y'-y$, $a\equiv\frac{k^+}{\vee}$, $c_0=2e^{-\gamma_E}$, $\hat\mu^2=\mu^2e^{2\gamma_E}/4=\mu^2/c_0^2$, $d_\perp=2-\epsilon$. $\frac{d^3N_{\xi\rm LO}}{d^3k}[f]$ is your (5.13) with the weight $f(X)$ \emph{inside} $\int_{x,y,x',y'}$; $\frac{d^3N_{\rm LO}}{d^3k}[f]$ likewise ($\xi=1$).
\begin{equation}
\frac{d^3N^y_{\rm LO}}{d^3k}\equiv\frac{1}{(2\pi)^3}\frac{g^2}{2\pi^2k^+}\int_{x,x',z}\frac{(x-z)\cdot(x'-z)}{(x-z)^2(x'-z)^2}\big[U^{e'c}(z)-U^{e'c}(x')\big]\big[U^{ec}(z)-U^{ce}(x)\big]\rho^{e'}(x')\rho^{e}(x),\label{eq:LOy}\end{equation}\begin{equation}H_4=\frac{g^2N_c}{\pi\mathbf k^2}\frac{d^3N^y_{\rm LO}}{d^3k}.
\end{equation}

\section{Step by step}
\paragraph{Step 0: the full Row IV} (your three lines, before anything is extracted; $\frac{\xi}{(1-\xi)_+}$ is your $p^+$ prescription, $p^+=k^+\frac{1-\xi}{\xi}$):
\begin{align}
\text{(L1)}&=\frac{g^2N_c}{4\pi^2}\int_a^1\frac{d\xi}{\xi^2}\Big[\xi(1-\xi)+\frac{1-\xi}{\xi}+\frac{\xi}{(1-\xi)_+}\Big]\frac{d^3N_{\xi\rm LO}}{d^3k}[\mathcal B_0],\qquad
\text{(L2)}=\frac{g^2N_c}{4\pi^2}\Big[\log\frac{k^+}{\Lambda}+\log\Big(\frac{\vee}{k^+}-1\Big)\Big]\frac{d^3N_{\rm LO}}{d^3k}[\mathcal B_0],\notag\\
\text{(L3)}&=\frac{g^2N_c}{4\pi^2}\int_a^{1}d\xi\,\frac{1-\xi}{\xi}\,\frac{d^3N^{\rm lp}_{\xi\rm LO}}{d^3k}.
\end{align}
Every term is symmetric under amplitude $\leftrightarrow$ conjugate. There is no $H_4$ and no $d^3N^y$.

\paragraph{Step 1: the $\epsilon$-bracket.} $d_\perp=2-\epsilon$ gives exactly $\mathcal B_0=-\frac2\epsilon-\frac{d_\perp}{2}\ln(X^2\hat\mu^2)\equiv\mathcal B$. In terms without a pole in front, $\frac{d_\perp}{2}\to1$.

\paragraph{Step 2: your $p^+$ prescription $\to$ the standard plus distribution; (L2) cancels except $\log\frac{k^+}{\Lambda}$.} Exactly (\texttt{g2iv\_full/plus\_check.py}, $6\times10^{-17}$):
\begin{equation}
\text{(L1)}=\frac{g^2}{8\pi^2}\int_a^1\frac{d\xi}{\xi^2}\Big[\mathcal P_{gg}-\frac{\beta_0}{2}\delta(1-\xi)\Big]\frac{d^3N_{\xi\rm LO}}{d^3k}[\mathcal B]-\frac{g^2N_c}{4\pi^2}\log\Big(\frac{\vee}{k^+}-1\Big)\frac{d^3N_{\rm LO}}{d^3k}[\mathcal B],
\end{equation}
where $\mathcal P_{gg}$ is the standard splitting function, with the standard plus distribution on $[a,1]$ (endpoint $\log(1-a)$ included). Hence
\begin{equation}
\text{(L1)}+\text{(L2)}=\underbrace{\frac{g^2N_c}{4\pi^2}\log\frac{k^+}{\Lambda}\frac{d^3N_{\rm LO}}{d^3k}[\mathcal B]}_{\rm JIMWLK}+\frac{g^2}{8\pi^2}\int_a^1\frac{d\xi}{\xi^2}\Big[\mathcal P_{gg}-\frac{\beta_0}{2}\delta(1-\xi)\Big]\frac{d^3N_{\xi\rm LO}}{d^3k}[\mathcal B].\label{eq:s2}
\end{equation}

\paragraph{Step 3: DGLAP.} The $\mathcal P_{gg}$ term is
\begin{equation}
{\rm (D)}\equiv-\frac{g^2}{8\pi^2}\int_a^1\frac{d\xi}{\xi^2}\,\mathcal P_{gg}(\xi)\,\frac{d^3N_{\xi\rm LO}}{d^3k}\Big[\frac2\epsilon+\ln(X^2\hat\mu^2)\Big].
\end{equation}
Its pole is the $\MSb$ counterterm of $D_{h/g}$: the region $a\le\zeta\xi\le\zeta\le1$ is that of $\int_{k^+/\vee}^1\frac{d\zeta}{\zeta^3}D(\zeta)$. After it, $\ln(X^2\hat\mu^2)\to\ln(X^2\mu_F^2/c_0^2)$.

\paragraph{Step 4: running coupling (your advisor's formula).} The $\delta$ term of \eqref{eq:s2} is $+\frac{g^2\beta_0}{16\pi^2}\frac{d^3N_{\rm LO}}{d^3k}\big[\frac2\epsilon+\ln(X^2\hat\mu^2)\big]$. The coupling counterterm of LO removes $\frac2\epsilon$, and $\frac{g^2}{16\pi^2}=\frac{\alpha_s}{4\pi}$, $\LO\propto\alpha_s$:
\begin{equation}
\frac{d^3N_{\rm LO}}{d^3k}\Big|_{\alpha_s}+\frac{\alpha_s\beta_0}{4\pi}\frac{d^3N_{\rm LO}}{d^3k}\big[\ln(X^2\hat\mu^2)\big]=\frac{d^3N_{\rm LO}}{d^3k}\Big|_{\alpha_s\to\alpha_s(X^2)},\qquad\alpha_s(X^2)=\alpha_s+\frac{\alpha_s^2\beta_0}{4\pi}\ln(X^2\hat\mu^2)
\end{equation}
(sympy: difference $0$, \texttt{advisor/checks.py}).

\paragraph{Step 5: the logarithm hidden in (D).}\label{sec:small} Look at (D) at small $\xi$. There $\mathcal P_{gg}\to2N_c\frac{1-\xi}{\xi}$, and the $y'$ integral is a Fourier transform at the large momentum $\mathbf k/\xi$:
\begin{equation}
I(\xi)\equiv\int d^2y'\,e^{-ik\cdot(y'-y)/\xi}\,\ln(X^2\hat\mu^2)\,F(y')=-\frac{4\pi\xi^2}{\mathbf k^2}F(y)\Big[1+O(\xi)\Big],\label{eq:I}
\end{equation}
where $F(y')$ is the smooth rest of the integrand (kernel, Wilson lines, charges, after $\int_{x'}$).
\begin{itemize}
\item The constant and $\ln\hat\mu^2$ multiply $(2\pi\xi)^2\delta^2(\mathbf k)=0$.
\item $\nabla^2\ln X^2=4\pi\delta^2(X)$ gives $-4\pi\xi^2/\mathbf k^2$.
\item At $y'=y$: $\frac{(x'-y')}{(x'-y')^2}\to\frac{(x'-y)}{(x'-y)^2}$ and $U(y')\to U(y)$, i.e.\ $F(y)$ is the integrand of $\frac{d^3N^y_{\rm LO}}{d^3k}$, with both gluon points at one point ($y\equiv z$).
\end{itemize}
\textbf{Numerical check} (\texttt{smallxi.py}, smooth model $F$, $|\mathbf k|=2$): $I(\xi)\big/\big(-\frac{4\pi\xi^2}{\mathbf k^2}F(y)\big)$:
\begin{center}\small
\setlength{\tabcolsep}{3pt}\begin{tabular}{@{}cccccc@{}}\toprule
$\xi$ & 0.5 & 0.2 & 0.1 & 0.05 & 0.03\\\midrule
ratio & $1.196-0.712i$ & $1.008-0.215i$ & $1.002-0.104i$ & $0.99994-0.048i$ & $0.99979-0.034i$\\
predicted $O(\xi)$ ($-i\xi$ here) & $-0.5i$ & $-0.2i$ & $-0.1i$ & $-0.05i$ & $-0.03i$\\\bottomrule
\end{tabular}
\end{center}
So, with $\frac{d\xi}{\xi^2}\cdot\frac{1-\xi}{\xi}\cdot\xi^2=\frac{(1-\xi)d\xi}{\xi}$, the small-$\xi$ end of (D) is
\begin{equation}
-\frac{g^2}{8\pi^2}\int_a^{1}\frac{d\xi}{\xi^2}\,2N_c\frac{1-\xi}{\xi}\Big(-\frac{4\pi\xi^2}{\mathbf k^2}\Big)\frac{d^3N^y_{\rm LO}}{d^3k}=\frac{g^2N_c}{\pi\mathbf k^2}\frac{d^3N^y_{\rm LO}}{d^3k}\Big[\log\frac{\vee}{k^+}-(1-a)\Big]=\Big[\log\frac{\vee}{k^+}-\Big(1-\frac{k^+}{\vee}\Big)\Big]H_4 .\label{eq:hidden}
\end{equation}
\textbf{This is where $+\log\frac{\vee}{k^+}H_4$ comes from}: the parent gluon with very large $p^+$ (BFKL region) inside the DGLAP integral.

\paragraph{Step 6: $H_4$ is the same object.} Your $H_4$ line has $\int_{w}e^{ik\cdot w}\frac{(w-z)^i}{(w-z)^2}=\frac{2\pi ik^i}{\mathbf k^2}e^{ik\cdot z}$ (checked), so $\int_{w,w'}(\dots)=\frac{4\pi^2}{\mathbf k^2}$ and $H_4=\frac{g^2N_c}{\pi\mathbf k^2}\frac{d^3N^y_{\rm LO}}{d^3k}$, Eq.~\eqref{eq:LOy}. Expanding its colour factor, $\log\frac{\vee}{k^+}H_4$ is exactly your four-line BFKL term (prefactor $2N_c\frac{g^4}{16\pi^5}=\frac{g^4N_c}{8\pi^5}$).

\paragraph{Step 7: extraction (add and subtract).}\label{sec:book} Write (D) as
\begin{align}
{\rm (D)}={}&\underbrace{-\frac{g^2}{8\pi^2}\int_a^1\frac{d\xi}{\xi^2}\Big\{\mathcal P_{gg}\frac{d^3N_{\xi\rm LO}}{d^3k}\Big[\frac2\epsilon+\ln(X^2\hat\mu^2)\Big]\new{-2N_c\frac{1-\xi}{\xi}\Big(-\frac{4\pi\xi^2}{\mathbf k^2}\Big)\frac{d^3N^y_{\rm LO}}{d^3k}}\Big\}}_{{\rm (D_{sub})}:\ \text{DGLAP, BFKL region subtracted}}\notag\\
&+\underbrace{\log\frac{\vee}{k^+}H_4}_{\to\ \rm BFKL}\;-\Big(1-\frac{k^+}{\vee}\Big)H_4 .
\end{align}
\begin{itemize}
\item In (D$_{\rm sub}$) the integrand is integrable at $\xi\to0$, by Eq.~\eqref{eq:I} (difference $O(\xi)$ relative). So (D$_{\rm sub}$) has no $\log\frac{\vee}{k^+}$.
\item The subtraction has no pole and no $\mu$, so the DGLAP renormalization (Step 3) is unchanged.
\item Every term is symmetric.
\end{itemize}

\paragraph{Step 8: Row IV, final.}
\begin{equation}
\boxed{\begin{aligned}
\text{Row IV}={}&\underbrace{\frac{g^2N_c}{4\pi^2}\log\frac{k^+}{\Lambda}\frac{d^3N_{\rm LO}}{d^3k}[\mathcal B]}_{\rm JIMWLK}+\underbrace{\log\frac{\vee}{k^+}H_4}_{\rm BFKL}\\
&-\frac{g^2}{8\pi^2}\int_a^1\frac{d\xi}{\xi^2}\Big\{\mathcal P_{gg}\frac{d^3N_{\xi\rm LO}}{d^3k}\Big[\frac2\epsilon+\ln(X^2\hat\mu^2)\Big]+\frac{8\pi N_c\,\xi(1-\xi)}{\mathbf k^2}\,\frac{d^3N^y_{\rm LO}}{d^3k}\Big\};\text{(DGLAP$_{\rm sub}$)}\\
&+\frac{g^2\beta_0}{16\pi^2}\frac{d^3N_{\rm LO}}{d^3k}\Big[\frac2\epsilon+\ln(X^2\hat\mu^2)\Big]\quad\text{(running coupling)}\\
&-\Big(1-\frac{k^+}{\vee}\Big)H_4+\frac{g^2N_c}{4\pi^2}\int_a^1d\xi\,\frac{1-\xi}{\xi}\frac{d^3N^{\rm lp}_{\xi\rm LO}}{d^3k}\Big|_{d_\perp=2}\quad\text{(finite)}
\end{aligned}}\label{eq:final}
\end{equation}
(the second term in the braces is Eq.~\eqref{eq:repl}: $-2N_c\frac{1-\xi}{\xi}\cdot(-\frac{4\pi\xi^2}{\mathbf k^2})=\frac{8\pi N_c\xi(1-\xi)}{\mathbf k^2}$).
\begin{itemize}
\item \textbf{Check:} adding the BFKL term back, the three $H_4$ structures cancel:
\begin{equation*}
\log\tfrac{\vee}{k^+}H_4\;\underbrace{-\big[\log\tfrac{\vee}{k^+}-(1-a)\big]H_4}_{\text{subtraction}}\;-(1-a)H_4=0 ,
\end{equation*}
and \eqref{eq:final} returns to Steps 0--4. This is your advisor's ``it should cancel with $H_4$''.
\item \textbf{Correction to Eq.~(5.12)}: replace
\begin{equation*}
\old{-2N_c\frac{1-\xi}{\xi}\frac{d^3N^y_{\xi\rm LO}}{d^3k}\Big[\log\frac{(y'-y)^2\mu^2}{4}\Big]}\quad\text{by}\quad\new{+\frac{8\pi N_c\,\xi(1-\xi)}{\mathbf k^2}\frac{d^3N^y_{\rm LO}}{d^3k}},
\end{equation*}
and drop (5.14). The value is the same, but it is now well defined and symmetric. Also put the brackets of the DGLAP and RC lines inside the integrals.
\item At leading power ($a\to0$): $\int_a^1\to\int_0^1$, $(1-a)H_4\to H_4$.
\end{itemize}

\section{Checks}
\begin{center}\small
\begin{tabular}{@{}p{0.62\textwidth}p{0.33\textwidth}@{}}\toprule
$p^+$ prescription $=$ standard plus $-F(1)\log(\frac{\vee}{k^+}-1)$; the cut $p^+\ge\Lambda$ $=$ your (L2) & $6\times10^{-17}$, $4\times10^{-16}$ (\texttt{g2iv\_full/plus\_check.py})\\
RC $=\alpha_s(X^2)$ (advisor's formula) & exact (\texttt{advisor/checks.py})\\
small-$\xi$ limit \eqref{eq:I}, ratio $\to1$, next order $-i\xi$ & table above (\texttt{smallxi.py})\\
$\int_w e^{ik\cdot w}\frac{(w-z)^i}{(w-z)^2}=\frac{2\pi ik^i}{\mathbf k^2}e^{ik\cdot z}$; $\int d^2r\,e^{-iq\cdot r}\ln r^2=-\frac{4\pi}{q^2}$ & $4.83322$ vs $4.83322$; $-7.45\to-7.44$\\
$H_4=\frac{g^2N_c}{\pi\mathbf k^2}\frac{d^3N^y_{\rm LO}}{d^3k}$; $\int_a^1\frac{1-\xi}{\xi}=\log\frac1a-(1-a)$ & exact\\\bottomrule
\end{tabular}
\end{center}
\end{document}
